% \iffalse meta-comment
%
% hit-thesis-float.dtx 是浮动体与图表
%
% \fi
%
% \section{浮动体与图表}
%
% ^^A ======================================================================
% ^^A Module thesis-floats: 浮动体默认行为与双语题注（hit-thesis）
% ^^A ======================================================================
%    \begin{macrocode}
%<*thesis-floats>
%    \end{macrocode}
%
% \changes{v3.0.0}{2020/05/21}{添加博后的浮动间距设置}
% \changes{v3.2a}{2026/08/08}{浮动体间距的三个分支迁到 expl3}
% 博后单独一套间距；其余按 \option{style/glue} 决定要不要给伸缩量。
%    \begin{macrocode}
\bool_if:NTF \g__hit_postdoc_bool
  {
    \skip_set:Nn \intextsep    { 10.5bp }
    \skip_set:Nn \textfloatsep { 10.5bp }
    \skip_set:Nn \floatsep     { 10.5bp }
  }
  {
    \bool_if:NTF \g__hit_glue_bool
      {
        \skip_set:Nn \intextsep    { 8.50398bp~plus~2.83465bp~minus~0bp }
        \skip_set:Nn \textfloatsep { 8.50398bp~plus~2.83465bp~minus~0bp }
        \skip_set:Nn \floatsep     { 12bp~plus~2.83465bp~minus~0bp }
      }
      {
        \skip_set:Nn \intextsep    { 8.50398bp }
        \skip_set:Nn \textfloatsep { 8.50398bp }
        \skip_set:Nn \floatsep     { 12bp }
      }
  }
%    \end{macrocode}
%
% 此处设置float在p选项时间隔，此处不设置\cs{@fptop}和\cs{@fpbot}以确保居中。
% \changes{v1.0.12}{2018/04/03}{修正float为p状态时默认不居中bug}
% \changes{v2.0.04}{2018/12/04}{删除\cs{@fpsep}设置，似乎没有什么用}
% \changes{v2.0.04}{2018/12/04}{更新\cs{intextsep}\cs{textfloatsep}\cs{floatsep}间距为正文行间距}
% 下面这组命令使浮动对象的缺省值稍微宽松一点，从而防止幅度对象占据过多的文本页面，
% 也可以防止在很大空白的浮动页上放置很小的图形。
% \changes{v1.0.08}{2017/11/05}{修改附录中图、表、公式数字编码}
%    \begin{macrocode}
\tl_gput_right:Nn \appendix { \cs_set_nopar:Npn \thefigure { \thechapter - \arabic {figure} } }
\tl_gput_right:Nn \appendix { \cs_set_nopar:Npn \thetable { \thechapter - \arabic {table} } }
\tl_gput_right:Nn \appendix { \cs_set_nopar:Npn \theequation { \thechapter - \arabic {equation} } }
\cs_set:Npn \textfraction      { 0.15 }
\cs_set:Npn \topfraction       { 0.85 }
\cs_set:Npn \bottomfraction    { 0.65 }
\cs_set:Npn \floatpagefraction { 0.60 }
%    \end{macrocode}
%
% 由于我工的双标题，导致标题之下多出一空白字符的距离，去除。
% \changes{v2.0.04}{2018/12/04}{更新图段后空白距离}
% \changes{v2.0.04}{2018/12/04}{删除表段后空白距离}
% \changes{v2.0.05}{2018/12/05}{删除图段后空白距离}
%    \begin{macro}{\@makecaption}
% 根据我工规范，本科和硕博的图题序号之后的空格不一样。
%    \begin{hitrgu}[\PGR][2.13.1]
% 每个图均应有图题（由图序和图名组成），图题不宜有标点符号，图名在图序之后空1个
% 半角字符排写。
%    \end{hitrgu}
%    \begin{hitrgu}[\UGR][2.13.1]
% 每个图均应有图题（由图序和图名组成），图题不宜有标点符号，图名在图序之后空1个
% 字符排写。
%    \end{hitrgu}
% 我工规范中没有明确规定是否标题是否居中对齐，这里给出一个居中选项自行调整。
% 注意，我工只规定：“居中书写”。此处不额外添加悬挂处理。
% \changes{v1.0.06}{2017/10/25}{此处更改了选项的名称}
% \changes{v1.0.07}{2017/11/04}{优化了最后一行居中算法，使其两边对齐、单词内部断行}
%
% \changes{v3.2a}{2026/08/07}{\cs{@makecaption} 与 \cs{LT@makecaption} 迁到 expl3}
% \changes{v3.2a}{2026/08/11}{题号与题名之间的间隔改成读
%   \option{after-figure-caption-number} 与 \option{after-table-caption-number}}
% \cs{hit@caption@separator} 是题号与题名之间的间隔。图题表题各读一个常量，
% 常量为空就退回原来的取值：本科一个汉字宽，其余是 \cs{enskip}。
% 图还是表看 \cs{@captype}；\cs{LT@makecaption} 那条路上它也是有的，取不到就按图算。
% \cs{hit@caption@last@line} 把最后一行单独取出来居中，\option{capcenterlast} 开时用。
% 里面的盒子编号与分组照抄原来的写法，只是换成 expl3 的拼法。量宽仍用 \cs{sbox}
% 而不是 \cs{hbox\_set:Nn}：前者会套一层 \cs{color@begingroup}，盒子里的颜色
% 改动不会漏出去，后者没有。
%    \begin{macrocode}
\cs_new:Npn \hit@caption@separator@fallback
  { \hit@if@option@TF { bachelor } { \hskip \ccwd } { \enskip } }
\cs_new:Npn \hit@caption@number@lang:
  { \bool_if:NTF \g__hit_en_bool { en } { zh } }
\cs_new:Npn \hit@caption@type:
  { \cs_if_exist:NTF \@captype { \@captype } { figure } }
\cs_new:Npn \hit@caption@separator
  {
    \str_if_eq:eeTF { \hit@caption@type: } { table }
      {
        \hit@after@number:nnn { table@caption } { \hit@caption@number@lang: }
          { \hit@caption@separator@fallback }
      }
      {
        \hit@after@number:nnn { figure@caption } { \hit@caption@number@lang: }
          { \hit@caption@separator@fallback }
      }
  }
\cs_new_protected:Npn \hit@caption@last@line #1
  {
    \vskip 6.3bp
    {
      \tex_setbox:D 0 = \tex_vbox:D {#1}
      \tex_setbox:D 1 = \tex_vbox:D
        {
          \tex_unvbox:D 0
          \tex_setbox:D 2 = \tex_lastbox:D
          \hbox_to_wd:nn { \textwidth }
            { \hfill \tex_unhcopy:D 2 \unskip \unskip \hfill }
        }
      \tex_unvbox:D 1
    }
  }
\cs_set:Npn \@makecaption #1#2
  {
    \vskip \abovecaptionskip
    \wuhao
    \sbox \@tempboxa { #1 \hit@caption@separator #2 }
    \dim_compare:nNnTF { \box_wd:N \@tempboxa } > { \hsize }
      {
        \hit@if@option@TF { capcenterlast }
          { \hit@caption@last@line { #1 \hit@caption@separator #2 } }
          { #1 \hit@caption@separator #2 }
        \par
      }
      {
        \tex_global:D \@minipagefalse
        \hb@xt@ \hsize { \hfil \box \@tempboxa \hfil }
      }
    \vskip \belowcaptionskip
  }
%    \end{macrocode}
%
%    \end{macro}
%    \begin{macro}{\LT@makecaption}
% 有关longtable的caption支持
%    \begin{macrocode}
\cs_set:Npn \LT@makecaption #1#2#3
  {
    \LT@mcol \LT@cols c
      {
        \hbox_to_zero:n
          {
            \hss
            \parbox [t] \LTcapwidth
              {
                \wuhao
                \sbox \@tempboxa { #1 { #2 \hit@caption@separator } #3 }
                \dim_compare:nNnTF { \box_wd:N \@tempboxa } > { \hsize }
                  {
                    \hit@if@option@TF { capcenterlast }
                      { \hit@caption@last@line { #1 { #2 \hit@caption@separator } #3 } }
                      { #1 { #2 \hit@caption@separator } #3 }
                    \par
                  }
                  {
                    \tex_global:D \@minipagefalse
                    \hbox_to_wd:nn { \hsize } { \hfil \box \@tempboxa \hfil }
                  }
                \endgraf \vskip \belowcaptionskip
              }
            \hss
          }
      }
  }
%    \end{macrocode}
%
%    \end{macro}
%    \begin{macro}{\longbionenumcaption}
% 长表格的双语标题是一个坑. 因为第一不能用浮动格式，只能用longtable包中的tabular
% ，这样表题只能使用表格中前两行来写。这样出现了一个问题是，中英表题的间距，标题
% 和表第一行间距，表格内部间距等多个变量的协调问题。这个问题只要使用tabular的形
% 式，就是无解的。唯一的方法就是把这些参数都给用户列出来。以下，第2，5参数为中英
% 双语标题内容，1，4为标题参数。6为中英标题间距，7为表题和表格间距。
%    \begin{macrocode}
\cs_set_nopar:Npn \longbionenumcaption #1 #2 #3 #4 #5 #6 #7
  {
    \noalign { \__hit_warn_once:nnn { longbionenumcaption } { longbicaption-deprecated } { } }
    \hit@long@bicaption {#1} {#2} {#3} {#4} {#5} {#6} {#7}
  }
\cs_set_nopar:Npn \hit@long@bicaption #1 #2 #3 #4 #5 #6 #7
  {
%
\@if@contemptyarg{#1}{\caption{#2}}{\caption[#1]{#2}}%
\cs_gset_eq:NN \@cont@oldtablename \tablename
\cs_gset_nopar:Npn \tablename {#3}~
\cs_gset_eq:NN \LT@c@ption \@cont@LT@nonumintoc
\\[#6]~
\@if@contemptyarg{#4}{\caption{#5}}{\caption[#4]{#5}}%
\cs_gset_eq:NN \tablename \@cont@oldtablename
\cs_gset_eq:NN \LT@c@ption \@cont@oldLT@c@ption
\vspace{#7}
  }
%    \end{macrocode}
%
%    \end{macro}
%
%    \begin{macro}{\LT@caption}
% \changes{v3.2a}{2026/08/10}{长表格的双语题注统一到 \cs{caption}\marg{中文}\oarg{英文}，
%   跟浮动体一致。\cs{longbionenumcaption} 照旧能用，上面那条就是靠它实现的}
% 浮动体里 \cs{caption}\marg{中文}\oarg{英文} 排双语，长表格里却只能写
% \cs{longbionenumcaption} 那七个参数。原因是 \env{longtable} 在 \cs{begin} 时把
% \cs{caption} 换成了自己的 \cs{LT@caption}，只认 \oarg{短标题}\marg{正文}，尾随的
% \oarg{英文} 会当成正文印出来。
%
% 这里给 \cs{LT@caption} 加一层：先自己接下 \texttt{*}、\oarg{短标题}、\marg{正文}，
% 再探一眼后面有没有 \oarg{英文}。没有就把接到的参数原样交回原命令——星号形式与
% \pkg{caption} 宏包那条路都不碰，行为一个字不变；有就转给
% \cs{longbionenumcaption}。
%
% 接参数与探看都在 \cs{noalign}\cs{bgroup} 里做，跟原命令一样：那个位置只做展开，
% \cs{futurelet} 这类赋值会当场开一个单元格，题注的 \cs{noalign} 就错位了。分叉之后
% 先 \cs{egroup} 退出来，再调原命令，它自己会重新进 \cs{noalign}。
%
% \cs{longbionenumcaption} 内部还要调 \cs{caption}，那时没有尾随可选参数，走的是
% 交回原命令那条，不会绕回来。
%    \begin{macrocode}
\cs_set_eq:NN \hit@original@LT@caption \LT@caption
\cs_set:Npn \LT@caption
  {
    \noalign \bgroup
      \@ifstar
        { \hit@lt@caption@star }
        { \@ifnextchar [ { \hit@lt@caption@short } { \hit@lt@caption@plain } }
  }
\cs_set:Npn \hit@lt@caption@star #1
  { \egroup \hit@original@LT@caption * {#1} }
\cs_set:Npn \hit@lt@caption@short [#1] #2
  { \hit@lt@caption@peek { [#1] } {#1} {#2} }
\cs_set:Npn \hit@lt@caption@plain #1
  { \hit@lt@caption@peek { } { } {#1} }
\cs_set:Npn \hit@lt@caption@peek #1#2#3
  {
    \@ifnextchar [
      { \hit@lt@caption@bilingual {#2} {#3} }
      { \egroup \hit@original@LT@caption #1 {#3} }
  }
\cs_set:Npn \hit@lt@caption@bilingual #1#2 [#3]
  {
    \egroup
    \hit@long@bicaption {#1} {#2} { \hit@table@prefix@en } { } {#3}
      { -0.5em } { 3.15bp }
  }
%    \end{macrocode}
%    \end{macro}

%    \begin{macro}{\ltfontsize}
% \changes{v3.2a}{2026/08/07}{长表格字号机制抽到 \file{src/utils/hit-table.dtx}}
% 我们采用 \pkg{longtable} 来处理跨页的表格。默认字号压到五号，机制见
% \file{src/utils/hit-table.dtx}。
%    \begin{macrocode}
\hit@setup@table
%    \end{macrocode}
%
%    \end{macro}
% 图表名称及格式。
% \changes{v1.0.08}{2017/11/05}{删除冗余公式符号定义}
%    \begin{macrocode}
\cs_set:Npn \thesubtable { ( \alph { subtable } ) }
\cs_set:Npn \thefigure   { \arabic { chapter } - \arabic { figure } }% 使图编号为 7-1 的格式
%    \end{macrocode}
%
% \changes{v3.0.15}{2020/05/21}{根据窝工2021“四月革命”之精神，修改子图为双括号}
%    \begin{macrocode}
\cs_set:Npn \thesubfigure { ( \alph { subfigure } ) }% 子图编号为 a) 的格式
\cs_set:Npn \p@subfigure  { \thefigure \nobreakspace { } }% 子图引用为 7-1 a)
\cs_set:Npn \thetable     { \arabic { chapter } - \arabic { table } }% 表编号为 7-1
%    \end{macrocode}
%
% \changes{v3.2a}{2026/08/07}{\cs{setitemize}/\cs{setenumerate} 从 book-floats 迁出至 book-paragraph}
% 此处删除hang caption的设置
% \changes{v1.0.06}{2017/10/25}{删除caption hang 的默认设置，因为不在规范要求中}
%    \begin{macrocode}
\captionnamefont{\wuhao}
\captiontitlefont{\wuhao}
\cs_set:Npn \subcapsize { \wuhao }
%    \end{macrocode}
%
% \option{modern} 下子图题的字体、标签分隔与最后一行对齐交给 \pkg{caption}，
% 键名不同但取的是同一批值。
%    \begin{macrocode}
\DeclareCaptionFont { hit@subcap } { \subcapsize }
\DeclareCaptionLabelSeparator { hit@sub } { \hskip \subfiglabelskip }
\captionsetup [ sub ]
  { font = hit@subcap , labelformat = simple , labelsep = hit@sub , skip = 0pt }
\bool_if:NT \g__hit_subcapcenterlast_bool
  { \captionsetup [ sub ] { justification = centerlast } }
\dim_set:Nn \abovecaptionskip {0pt}%为了双标题之间的间距，不能设置
\dim_set:Nn \belowcaptionskip {0pt}
%    \end{macrocode}
%
% \changes{v3.2a}{2026/08/10}{\env{publist} 挪到 \file{src/utils/hit-list.dtx}，各个类共用}
% \env{publist} 见 \file{src/utils/hit-publist.dtx}。
%    \begin{macrocode}
\hit@setup@publist
%    \end{macrocode}
%
% \changes{v2.0.01}{2018/06/28}{去除定理注释括号}
% \changes{v3.2a}{2026/08/07}{ntheorem 定理样式与 \cs{arraycolsep} 从 book-floats 迁出至 book-math；\cs{setlength}\{\cs{parindent}\}\{2em\} 迁出至 book-paragraph}
%    \begin{macrocode}
%</thesis-floats>
%    \end{macrocode}
%
% \Finale
